% ════════════
% wp7_manuscript.tex
% Material Dignity Infrastructure: Prototype Capital Architecture and First-Mover Risk Specification
% Material Dignity Infrastructure Series
% ════════════

\documentclass[12pt, letterpaper]{article}

% ── PACKAGES ──
\usepackage[left=0.7in, right=0.7in, top=1in, bottom=1in, headheight=14pt]{geometry}
\usepackage[expansion=false]{microtype}
\usepackage{textcomp}
\usepackage{setspace}
\usepackage{booktabs}
\usepackage{tabularx}
\usepackage[titles]{tocloft}
\usepackage{tabularray}
\usepackage{longtable}
\usepackage{array}
\usepackage{multirow}
\usepackage{hyperref}
% natbib not used — manual thebibliography
\usepackage{amsmath}
\usepackage{amssymb}
\usepackage{parskip}
\usepackage{titlesec}
\usepackage{fancyhdr}
\usepackage[table]{xcolor}
\usepackage{caption}
\usepackage{footnote}
\usepackage{enumitem}
\usepackage{tikz}
\usepackage{needspace}
\usepackage{etoolbox}
\usepackage{float}

% ── COLOR SYSTEM ───
\definecolor{auditblue}{HTML}{003366}
\definecolor{rowgray}{HTML}{F5F5F5}

% ── GLOBAL SPACING ───
\setstretch{1.4}
\setlength{\parindent}{0pt}
\setlength{\parskip}{8pt}
\hyphenpenalty=1000

% ── SECTION FORMATTING. ZERO BOLDING IN TEXT. ───
\titleformat{\section}
  {\normalfont\large\color{auditblue}}{\thesection.}{0.5em}{}
\titleformat{\subsection}
  {\normalfont\normalsize\color{auditblue}}{\thesubsection.}{0.5em}{}
\titlespacing*{\section}{0pt}{18pt}{6pt}
\titlespacing*{\subsection}{0pt}{12pt}{4pt}

% ── HEADER / FOOTER ───
\pagestyle{fancy}
\fancyhf{}
\rhead{\small\textit{{Material Dignity Infrastructure Series}}}
\lhead{\small\textit{Working Paper, June 2026}}
\cfoot{\thepage}
\renewcommand{\headrulewidth}{0.4pt}

% ── TABLE SETTINGS ───
\renewcommand{\arraystretch}{1.3}

% ── ORPHAN / WIDOW CONTROL ───
\clubpenalty=10000
\widowpenalty=10000
\displaywidowpenalty=10000

% ── BIBLIOGRAPHY ORPHAN PREVENTION ───
\AtBeginEnvironment{thebibliography}{\interlinepenalty=10000}

% ── HYPERREF CONFIGURATION ───
\hypersetup{
  colorlinks=true,
  linkcolor=black,
  citecolor=black,
  urlcolor=black,
  pdftitle={{Material Dignity Infrastructure}: {Prototype Capital Architecture and First-Mover Risk Specification}},
  pdfauthor={{Charles J. DiBella}},
  pdfsubject={Working Paper},
  pdfkeywords={{Permanent Supportive Housing, Catalytic Equity, CalAIM, Blended Finance, Real Estate Acquisition, Homelessness Policy, Avoided Cost Returns, Community Supports}}
}
\setcounter{tocdepth}{2}

\begin{document}
\captionsetup[table]{labelfont=normalfont, labelsep=period, skip=6pt}

% ── FRONT MATTER ───
\begin{titlepage}
\begin{center}
\setstretch{1.4}
{\LARGE \textbf{Material Dignity Infrastructure}}\\[6pt]
{\large Prototype Capital Architecture and First-Mover Risk Specification}\\[0.5cm]
{\normalsize \textbf{Charles J. DiBella}}\\[0.01cm]
{\normalsize Director of Urban Systems Research}\\[0.01cm]
{\normalsize Working Paper, June 2026}\\[0.3cm]
\textbf{ABSTRACT}
\vspace{0.2cm}
\end{center}
\begin{minipage}{0.95\textwidth}
\setstretch{1.15}
\noindent

The preceding six papers in this series established the governance architecture, diagnosed national structural misalignment, specified the Los Angeles metropolitan stabilization pipeline, defined the relational layer required within that pipeline, operationalized the field manuals for autonomic routing and social network transition, and cataloged the political economy of opposition. Each paper assumed a prototype would be built. None specified who builds it, through what financial instrument, or how first-mover capital risk is absorbed before the self-executing replication mechanism activates. This paper closes that gap. It translates the Material Dignity Infrastructure into the language capital markets require before committing funds: a capitalization table with defined tranches, a revenue model anchored to sovereign-backed Medi-Cal billing codes, a thirty-six-month occupancy ramp projecting stabilized yield, and a governance structure that insulates physical assets from operational liability. The result is an investable asset specification for the acquisition, conversion, and operation of One California Plaza as the MDI prototype site.

\end{minipage}
\end{titlepage}

\pagenumbering{arabic}

% ── DOCUMENT METADATA ────
\newpage
\section*{Document Metadata}

\noindent\textbf{Keywords}\\[1pt]
Permanent Supportive Housing, Catalytic Equity, CalAIM, Blended Finance, Real Estate Acquisition, Homelessness Policy, Avoided Cost Returns, Community Supports

\vspace{8pt}
\noindent\textbf{JEL Classification}
\begin{itemize}[nosep, before={\vspace{-8pt}}, leftmargin=2em]
    \item I18 Health: Government Policy • Regulation • Public Health
    \item I38 Welfare, Well-Being, and Poverty: Government Programs • Provision and Effects of Welfare Programs
    \item R38 Real Estate Markets, Spatial Production Analysis, and Firm Location: Government Policy • Regulatory Policies
    \item H43 Project Evaluation • Social Discount Rate
\end{itemize}

\vspace{8pt}
\noindent\textbf{Target eJournals}
\begin{itemize}[nosep, before={\vspace{-8pt}}, leftmargin=2em]
    \item Urban Economics eJournal
    \item Health Economics Network
    \item Public Economics: Expenditure \& Procurement eJournal
    \item Real Estate eJournal
\end{itemize}

\vspace{8pt}
\noindent\textbf{Suggested Citation}\\[1pt]
DiBella, C.J. (2026). Material Dignity Infrastructure: Prototype Capital Architecture and First-Mover Risk Specification. Material Dignity Infrastructure Series, June 2026.

\newpage

\vspace{8pt}
\noindent\textbf{Statement of Necessity}\\[1pt]
This paper provides the required financial specification to convert theoretical homelessness resolution frameworks into investable real estate assets, directly solving the first-mover capital risk gap.

\vspace{8pt}
\noindent\textbf{Author's Note}\\[1pt]
This paper is the fifth working paper in the Material Dignity Infrastructure series. It supplies the terminal financial specification required to capitalize the prototype.

% ── TABLE OF CONTENTS ────
\newpage
\begingroup
\setstretch{1.35}
\setlength{\cftbeforesecskip}{2pt}
\setlength{\cftbeforesubsecskip}{-2pt}
\tableofcontents
\endgroup
% ── BODY ────
\clearpage
\setstretch{1.35}

\clearpage
\section{The First-Mover Problem}

Every self-executing replication mechanism requires a demonstration that precedes replication. Capital markets do not replicate theoretical models. They replicate observed financial performance. The MDI architecture specifies that once One California Plaza demonstrates per capita cost reductions exceeding fifty thousand dollars at the three-hundred-sixty-five-day measurement window, institutional capital replicates the mechanism autonomously through standard real estate investment channels. That replication sequence is credible because it relies on proven capital market behavior rather than philanthropic goodwill or legislative appropriation.

The sequence contains a structural vulnerability at origin. The prototype must produce the financial results before those results can trigger replication. Someone must fund the demonstration before the demonstration funds everything that follows. That first actor cannot be responding to proven financial viability because the proof does not yet exist. The first actor accepts demonstration risk in exchange for a structured return that compensates for the absence of operating history.

This is not a social movement problem. Social movements generate political pressure; they do not capitalize real estate acquisitions. This is not a government procurement problem. Municipal procurement channels are structurally captured by incumbent service providers whose institutional survival depends on preserving existing funding architectures. This is a venture capital problem with a specific solution: a blended capital stack that layers catalytic first-loss equity beneath concessionary senior debt, producing a risk-adjusted return profile acceptable to infrastructure-class investors while absorbing the volatility of an unproven operating model.

\clearpage
\section{Acquisition Vehicle: One California Plaza}

One California Plaza occupies 1.05 million rentable square feet in the Bunker Hill district of downtown Los Angeles. The building was constructed in 1985 as Class A office space. Commercial office vacancy rates in downtown Los Angeles exceeded thirty percent by late 2025, driven by remote work adoption, corporate footprint consolidation, and deferred capital expenditure cycles. Buildings of this vintage and classification trade at steep discounts beneath peak valuation and replacement cost.

The acquisition targets an entry price not exceeding one hundred fifty dollars per rentable square foot. At 1.05 million square feet, total acquisition cost targets one hundred fifty-seven million five hundred thousand dollars. This entry price generates an immediate capital cushion relative to replacement cost, which for comparable downtown high-rise construction exceeds six hundred dollars per square foot. The cushion absorbs specialized conversion costs without pushing total capitalization above market parity for newly constructed permanent supportive housing.

\textbf{Conversion Scope.} The physical conversion repurposes commercial floor plates into clustered residential pods. Each pod houses four to six individuals in private rooms organized around shared communal space. Acoustic insulation, high-durability fixtures, secured storage, and pet accommodation constitute the primary retrofit requirements. Commercial HVAC, electrical, and plumbing infrastructure requires adaptation rather than replacement. Estimated conversion cost ranges from forty to sixty dollars per square foot, totaling forty-two million to sixty-three million dollars across the full building footprint.

\textbf{Total Capitalization.} Acquisition plus conversion produces a base capitalization requirement between one hundred ninety-nine million five hundred thousand and two hundred twenty million five hundred thousand dollars. Adding a capitalized interest reserve of ten million five hundred sixty thousand dollars to pre-fund twelve months of senior debt service during the ramp-up period brings total capitalization to two hundred twenty million dollars.

\clearpage
\section{Capitalization Table}

The capital stack layers two distinct tranches with structurally separated risk profiles.

\textbf{Tranche A: First-Loss Catalytic Equity (40\%).} Eighty-eight million dollars. Sourced through Program-Related Investments from philanthropic endowments, mission-driven family offices, and impact-first institutional allocators. This tranche absorbs one hundred percent of first-loss exposure during the ramp-up period. Return is subordinated to Tranche B and targets a nominal four percent annual yield, deferred until Month 24 and payable from operating surplus only after senior debt service is satisfied. The catalytic tranche accepts below-market return in exchange for structural proof that the MDI model converts distressed commercial real estate into stabilized permanent supportive housing at scale.

\textbf{Tranche B: Senior Concessionary Debt (60\%).} One hundred thirty-two million dollars. Sourced through Community Development Financial Institutions, New Markets Tax Credit allocations, and mission-aligned commercial lenders. Senior position in the capital stack with first claim on all operating revenue after direct operating expenses. Target yield of eight percent annual return. Debt service coverage ratio floors at 1.25x, meaning net operating income must exceed debt service obligations by twenty-five percent before any distributions flow to Tranche A equity holders.

\textbf{Structural Protection.} The Delaware Statutory Trust holds property title. The California Public Benefit Corporation operates the facility under a Master Lease Agreement. If the PBC defaults on lease obligations or enters insolvency, the lease reassignment clause automatically transfers operational authority to a pre-designated successor operator. The physical asset never enters bankruptcy proceedings. Tranche B debt holders retain their secured position against the real property regardless of operational performance.

\clearpage
\section{Revenue Model: CalAIM Billing Architecture}

The revenue model replaces commercial tenant rent with Medi-Cal managed care capitation. California's CalAIM initiative authorizes managed care plans to fund housing-related services through Enhanced Care Management and Community Supports billing codes. These codes provide the operational revenue stream that services senior debt and generates stabilized yield.

\textbf{Revenue Line 1: Housing Transition Navigation Services.} Billing code H0043 with modifier U6. Reimburses street-level outreach, trust accumulation, demographic enrollment, and pre-placement coordination. Per-member-per-month rate structure. Active during Phase Zero before physical placement. Estimated annual revenue per active outreach case: eight thousand four hundred dollars.

\textbf{Revenue Line 2: Housing Deposits.} Billing code H0044 with modifier U2. Reimburses upfront capital expenditures associated with unit preparation, security deposits, and initial furnishing. One-time billing event per placement. Estimated per-unit revenue: three thousand five hundred dollars.

\textbf{Revenue Line 3: Housing Tenancy and Sustaining Services.} Billing code T2050 with modifier U6. Reimburses ongoing tenancy support, relational co-regulation, pod-steward coordination, and crisis de-escalation. Per-member-per-month rate structure. Active for twenty-four months post-placement. Estimated annual revenue per occupied unit: fourteen thousand four hundred dollars.

\textbf{Revenue Line 4: Recuperative Care.} Billing code T2033 with modifier U6. Reimburses medical respite services for individuals presenting acute metabolic instability, severe withdrawal, or sustained dorsal-vagal shutdown exceeding seventy-two hours. Higher per-diem rate than standard tenancy support. Isolated billing channel prevents acute stabilization costs from depleting the primary housing revenue pool. Estimated annual revenue per recuperative care bed: twenty-one thousand nine hundred dollars.

\textbf{Blended Revenue Per Occupied Unit.} At stabilized occupancy, the weighted average annual revenue per occupied unit across all four billing lines equals approximately eighteen thousand two hundred dollars. This figure excludes Enhanced Care Management capitation, which layers additional per-member-per-month revenue for individuals enrolled in the highest-acuity population tiers.

\clearpage
\section{Thirty-Six-Month Occupancy and Revenue Projection}

The occupancy ramp reflects the biological reality that trust accumulation and autonomic stabilization require time. Filling units faster than relational protocols permit degrades retention outcomes and produces the same revolving-door failure that characterizes existing shelter models.

\textbf{Months 1 through 6 (Ramp-up).} Target occupancy: thirty percent of converted capacity. Priority intake channels Recuperative Care activations, capturing individuals presenting acute physiological instability. Revenue composition skews toward higher-rate T2033 billing. Estimated gross revenue: thirteen million one hundred forty thousand dollars annualized.

\textbf{Months 7 through 18 (Stabilization).} Target occupancy: seventy-five percent of converted capacity. Billing composition transitions toward standard Housing Tenancy and Sustaining Services as acute cases stabilize and step down to community pods. Estimated gross revenue: twenty-eight million three hundred fifty thousand dollars annualized.

\textbf{Months 19 through 36 (Maturity).} Target occupancy: ninety-five percent of converted capacity. Revenue composition fully stabilized across all four billing lines. Estimated gross revenue: thirty-five million nine hundred ten thousand dollars annualized.

\textbf{Operating Expenses.} Direct operating costs including staffing, facilities maintenance, insurance, and administrative overhead are estimated at sixty percent of gross revenue at maturity, producing net operating income of thirteen million eight hundred thirty-two thousand dollars annually.

\textbf{Debt Service.} Senior debt service on one hundred thirty-two million dollars at eight percent annual yield requires ten million five hundred sixty thousand dollars annually. Net operating income at maturity covers senior debt service at a 1.31x coverage ratio, exceeding the 1.25x floor by six percentage points.

\textbf{Equity Return.} Residual operating surplus after senior debt service equals three million two hundred seventy-two thousand dollars annually, producing a nominal return of 3.7 percent on eighty-eight million dollars of catalytic equity. This falls marginally below the four percent target yield; layering Enhanced Care Management capitation revenue, excluded from the base model as conservative upside, closes the gap.

\clearpage
\section{Avoided-Cost Return Calculation}

The financial case for replication rests on a parallel calculation that does not appear on the facility's income statement but dominates the public fiscal argument. Los Angeles County spends approximately fifty thousand dollars per chronically unsheltered individual annually through emergency department admissions, psychiatric holds, law enforcement contacts, and sanitation operations. Permanent supportive housing targeted to intensive service users generates gross public savings exceeding forty-six thousand dollars per person per year by eliminating these recurring expenditures.

At stabilized occupancy of approximately two thousand individuals, One California Plaza generates estimated annual avoided public costs of one hundred million dollars. This figure represents the fiscal argument that triggers autonomous replication. When a county supervisor, a managed care CEO, or a commercial real estate investor observes that a single converted office tower eliminates one hundred million dollars in annual public waste while generating eight percent yield on senior debt secured by real property, the replication decision ceases to be philanthropic. It becomes fiduciary.

\clearpage
\section{Measurement Windows and Replication Triggers}

The prototype operates under three measurement windows that produce the evidentiary basis for replication.

\textbf{Day 90.} Retention rate measured against initial placement cohort. Target: eighty-five percent. This metric validates the autonomic routing matrix and social network transition protocol. If individuals placed through relational intake with preserved survival networks retain housing at rates exceeding eighty-five percent within the first ninety days, the biological intake thesis is operationally confirmed.

\textbf{Day 180.} Blended per-unit revenue measured against pro forma projection. Target: within ten percent of projected revenue per occupied unit. This metric validates the CalAIM billing architecture. If managed care reimbursement flows at projected rates, the revenue model is confirmed as operationally viable.

\textbf{Day 365.} Per capita avoided-cost reduction measured against baseline emergency service utilization. Target: verified reduction exceeding fifty thousand dollars per individual. This metric constitutes the replication trigger. Once the avoided-cost calculation is verified through audited emergency service data, the financial case for acquiring and converting the next building becomes self-evident to institutional capital.

\clearpage
\section{Risk Architecture}

\textbf{Occupancy Risk.} The ramp-up period between Month 1 and Month 18 generates revenue below stabilized projections. First-loss catalytic equity absorbs this shortfall by deferring return until Month 24. Senior debt service during the ramp-up period is funded through capitalized interest reserves built into the initial capital raise.

\textbf{Regulatory Risk.} CalAIM billing codes are authorized under federal Medicaid waiver authority. Waiver renewal cycles introduce regulatory uncertainty. Mitigation: the facility's physical asset retains intrinsic real estate value independent of operating revenue. If CalAIM authorization were revoked, the DST could release the property for alternative use at values supported by the distressed acquisition discount.

\textbf{Operational Risk.} The PBC operator may underperform on retention metrics or fail to execute relational protocols at specification. Mitigation: the Master Lease reassignment clause transfers operations to a successor entity without disrupting the capital structure. Senior debt holders retain their secured claim against the physical asset.

\textbf{Political Risk.} Federal policy shifts, state budget contractions, or municipal political opposition may impede operations. Mitigation: the tripartite structure deliberately bypasses municipal procurement channels. Revenue flows directly between managed care organizations and the PBC operator. No general fund appropriation is required. Political opposition cannot defund a revenue stream it does not control.

\clearpage
\section{Conclusion}

The Material Dignity Infrastructure series has produced a complete theoretical and operational architecture for resolving chronic unsheltered homelessness through neurologically calibrated, relationally engineered permanent supportive housing. This paper supplies the missing terminal specification: the financial instrument that makes the prototype investable. A blended capital stack of eighty-eight million dollars in catalytic first-loss equity and one hundred thirty-two million dollars in senior concessionary debt acquires and converts One California Plaza into a two-thousand-unit permanent supportive housing facility generating stabilized eight percent yield on senior debt, secured by real property held in an independent trust, and funded through sovereign-backed Medi-Cal billing codes that replace commercial tenant rent.

The prototype exists to be observed. Once the Day 365 measurement window validates per capita avoided-cost reductions exceeding fifty thousand dollars, the replication argument shifts from prospective to operational. Capital markets replicate proven financial performance. The first building must be funded by actors willing to accept demonstration risk. Every subsequent building funds itself.

\newpage

\addcontentsline{toc}{section}{References}
\section*{References}

\textit{   DiBella, C. J. (2025). }Material Dignity Infrastructure: A Distributed Stewardship Model for Homelessness Resolution*. SSRN Working Paper 5968756.
\textit{   DiBella, C. J. (2026). }Material Dignity Infrastructure: Structural Misalignment and the Activation of Surplus Shelter Capacity*. SSRN Working Paper 6211658.
\textit{   DiBella, C. J. (2026). }Material Dignity Infrastructure: Los Angeles Metropolitan Stabilization --- A Street-to-Home Pipeline Analysis*. SSRN Working Paper 6579600.
\textit{   DiBella, C. J. (2026). }Relational Dignity Infrastructure: The Human Layer Required to Make Material Housing Inhabitable*. SSRN Working Paper 6881539.
\textit{   DiBella, C. J. }Building Material Dignity: An Urban Stabilization Blueprint for the Unhoused*. ASIN: B0H1LSHYMF.
\textit{   DiBella, C. J. }The Architecture of Survival: An Infrastructure Blueprint for an Age of Collapse*. ASIN: B0H24BFK9F.
\textit{   California State Auditor (2024). }Homelessness in California: The State Must Do More to Assess the Cost-Effectiveness of Its Programs*.
\textit{   Culhane, D. P., Metraux, S., and Hadley, T. (2002). Public Service Reductions Associated with Placement of Homeless Persons with Severe Mental Illness in Supportive Housing. }Housing Policy Debate*, 13(1), 107-163.
\textit{   Olson, M. (1965). }The Logic of Collective Action: Public Goods and the Theory of Groups*. Harvard University Press.
\textit{   Porges, S. W. (1994). Orienting in a Defensive World: Mammalian Modifications of our Evolutionary Heritage. A Polyvagal Theory. }Psychophysiology*, 32(4), 301-318.


\clearpage
\section{Appendix A: Pro Forma Financial Model}

\subsection{Capital Requirements}

| Line Item | Amount | Notes |
|:---|:---|:---|
| Acquisition (1.05M sq ft at \$150/sq ft) | \$157,500,000 | Class B/C distressed entry price |
| Conversion (\$50/sq ft midpoint) | \$52,500,000 | Acoustic insulation, fixtures, pod configuration |
| Capitalized Interest Reserve | \$10,080,000 | 12 months senior debt service pre-funded |
| \textbf{Total Capitalization} | \textbf{\$220,080,000} | Rounded to \$220M for modeling |

\subsection{Capital Stack}

| Tranche | Amount | Percentage | Source | Target Yield |
|:---|:---|:---|:---|:---|
| A: Catalytic Equity (First-Loss) | \$88,000,000 | 40\% | PRIs, mission-driven family offices | 4.0\% (deferred to Month 24) |
| B: Senior Concessionary Debt | \$132,000,000 | 60\% | CDFIs, NMTC, mission-aligned lenders | 8.0\% |
| \textbf{Total} | \textbf{\$220,000,000} | \textbf{100\%} | | |

\subsection{CalAIM Revenue Lines (Per Unit, Annualized)}

| Billing Code | Service | Annual Revenue Per Unit |
|:---|:---|:---|
| H0043 U6 | Housing Transition Navigation | \$8,400 |
| H0044 U2 | Housing Deposits | \$3,500 (one-time) |
| T2050 U6 | Housing Tenancy and Sustaining | \$14,400 |
| T2033 U6 | Recuperative Care | \$21,900 |
| \textbf{Blended Average (stabilized)} | | \textbf{\$18,200} |

\subsection{Thirty-Six-Month Revenue Projection}

| Period | Months | Occupancy | Units Occupied | Gross Annual Revenue | Revenue Composition |
|:---|:---|:---|:---|:---|:---|
| Ramp-up | 1-6 | 30\% | 600 | \$13,140,000 | 60\% Recuperative Care, 40\% Tenancy |
| Stabilization | 7-18 | 75\% | 1,500 | \$27,300,000 | 30\% Recuperative Care, 70\% Tenancy |
| Maturity | 19-36 | 95\% | 1,900 | \$34,580,000 | 15\% Recuperative Care, 85\% Tenancy |

\subsection{Operating Expense Model (At Maturity)}

| Category | Annual Cost | Percentage of Gross Revenue |
|:---|:---|:---|
| Staffing (pod stewards, clinical, admin) | \$12,132,000 | 35\% |
| Facilities (maintenance, utilities, insurance) | \$5,187,000 | 15\% |
| Administrative overhead | \$3,458,000 | 10\% |
| \textbf{Total Operating Expenses} | \textbf{\$20,748,000} | \textbf{60\%} |
| \textbf{Net Operating Income} | \textbf{\$13,832,000} | \textbf{40\%} |

\subsection{Debt Service and Yield Distribution (At Maturity)}

| Line Item | Annual Amount |
|:---|:---|
| Senior Debt Service (8.0\% on \$132M) | \$10,560,000 |
| \textbf{Debt Service Coverage Ratio} | \textbf{1.31x} |
| Residual After Debt Service | \$3,272,000 |
| Catalytic Equity Return (on \$88M) | 3.7\% nominal |

\subsection{Avoided-Cost Fiscal Impact}

| Metric | Value |
|:---|:---|
| Annual street management cost per individual | \$50,000 |
| Individuals housed at maturity | 1,900 |
| Gross annual avoided public cost | \$95,000,000 |
| Net avoided cost (minus facility operating subsidy) | \$74,252,000 |

\subsection{Sensitivity Analysis}

| Variable | Base Case | Downside (-15\%) | Upside (+10\%) |
|:---|:---|:---|:---|
| Stabilized Occupancy | 95\% | 81\% | 100\% |
| Gross Revenue (Maturity) | \$34,580,000 | \$29,393,000 | \$36,400,000 |
| Net Operating Income | \$13,832,000 | \$11,757,200 | \$14,560,000 |
| DSCR | 1.31x | 1.11x | 1.38x |
| Equity Return | 3.7\% | 1.1\% | 4.5\% |

\subsection{Measurement Windows}

| Window | Day | Target Metric | Replication Trigger |
|:---|:---|:---|:---|
| Retention Validation | 90 | 85\% retention rate | Confirms autonomic routing protocol |
| Revenue Validation | 180 | Within 10\% of projected per-unit revenue | Confirms CalAIM billing architecture |
| Avoided-Cost Validation | 365 | >\$50,000 per capita reduction | \textbf{Triggers autonomous capital replication} |

\subsection{Key Assumptions and Limitations}

*   Unit count of 2,000 assumes full building conversion; partial conversion reduces proportionally.
*   CalAIM billing rates reflect 2025-2026 published rate structures and are subject to annual adjustment.
*   Recuperative Care revenue assumes dedicated floor allocation of approximately fifteen percent of total capacity.
*   Operating expense ratio of sixty percent is conservative relative to comparable permanent supportive housing operations.
*   Capitalized interest reserve covers twelve months of senior debt service during ramp-up, funded from initial capital raise.
*   Pro forma does not model Enhanced Care Management capitation overlay, which would increase per-unit revenue for highest-acuity enrollees.
*   Sensitivity analysis models a fifteen percent downside occupancy shortfall as the stress case.

\end{document}
